% Generado desde README.md. Editar libremente.

\begin{landscape}
\section{7. Datos sensibles y secretos del cliente}
No todos los datos valen lo mismo: el nombre de un bowl es público, pero el RUN, la dirección o el historial de pedidos de Javiera no lo son. Esta sección clasifica cada dato, decide quién puede verlo y cómo se protege, aclara que el sistema no guarda secretos de los clientes (ni datos de tarjeta) y termina con las pruebas que se ejecutaron contra el sistema funcionando.

\subsection{Clasificación de datos}
\vspace{0.6em}
{\scriptsize
\setlength{\LTpre}{0pt}\setlength{\LTpost}{0pt}
\setlength{\tabcolsep}{3pt}\renewcommand{\arraystretch}{1.15}
\begin{longtable}{|>{\raggedright\arraybackslash}p{\dimexpr 0.1441\linewidth-2\tabcolsep\relax}|>{\raggedright\arraybackslash}p{\dimexpr 0.1129\linewidth-2\tabcolsep\relax}|>{\raggedright\arraybackslash}p{\dimexpr 0.1061\linewidth-2\tabcolsep\relax}|>{\raggedright\arraybackslash}p{\dimexpr 0.0863\linewidth-2\tabcolsep\relax}|>{\raggedright\arraybackslash}p{\dimexpr 0.1239\linewidth-2\tabcolsep\relax}|>{\raggedright\arraybackslash}p{\dimexpr 0.0880\linewidth-2\tabcolsep\relax}|>{\raggedright\arraybackslash}p{\dimexpr 0.1706\linewidth-2\tabcolsep\relax}|>{\raggedright\arraybackslash}p{\dimexpr 0.1681\linewidth-2\tabcolsep\relax}|}
\hline
\textbf{Dato o credencial} & \textbf{Finalidad} & \textbf{Clasificación} & \textbf{Sistema custodio} & \textbf{Roles autorizados} & \textbf{En el frontend} & \textbf{Protección} & \textbf{Retención} \\ \hline
\endfirsthead
\hline
\textbf{Dato o credencial} & \textbf{Finalidad} & \textbf{Clasificación} & \textbf{Sistema custodio} & \textbf{Roles autorizados} & \textbf{En el frontend} & \textbf{Protección} & \textbf{Retención} \\ \hline
\endhead
Nombre, descripción, precio de productos & Vender & Público & inventario & Todos; editar admin/dueño & Menú & Validación en servidor & Mientras exista el producto \\ \hline
Stock & Evitar sobreventa & Interno & inventario & admin, dueño & Solo panel & DTO por rol & Siempre actualizado \\ \hline
RUN & Boleta e identificación del comprador (caso) & Personal & clientes (y copia en el pedido) & El cliente, admin, dueño & Mis datos, boleta & No sale a cocina ni despacho; no va en logs & Mientras la cuenta exista; en la boleta, lo que exija la ley tributaria \\ \hline
Nombre y teléfono & Contacto y entrega & Personal & clientes (copia en el pedido) & Cliente, personal & Panel, orden & Mínimo necesario & Igual \\ \hline
Dirección y comuna & Despacho & Personal & clientes (copia si es despacho) & Cliente, admin, cocina y despacho en el pedido & Orden de despacho & No va en el catálogo ni en logs & Igual \\ \hline
Correo & Cuenta, boleta & Personal & clientes y Auth & Cliente, admin & Mis datos & No va en logs (el notificador registra solo el dominio) & Igual \\ \hline
Fecha de nacimiento, sexo & Requeridos por el caso para el registro & Personal & clientes & Cliente, admin & Mis datos & Nunca salen en pedidos & Igual \\ \hline
Contraseña & Autenticar & Credencial & Auth (solo el hash) & Nadie la ve & Formulario & scrypt (N=16384) + sal por cuenta + pepper en Vault; nunca se devuelve & El hash, mientras exista la cuenta \\ \hline
Código de verificación & Verificar correo & Credencial temporal & Auth (solo su hash) & El dueño del correo & Correo simulado (modo demo) & SHA-256, vence en 24 h, se borra al usarlo & 24 h \\ \hline
Token de sesión & Mantener la sesión & Credencial & Auth (hash) y cookie & El navegador del usuario & Nunca lo ve JavaScript (HttpOnly) & 256 bits, opaco, SameSite, Path=/api & 30 min sin uso u 8 h \\ \hline
Pedido, líneas, totales & Vender y preparar & Personal-comercial & pedidos & Dueño del pedido, personal según rol & Orden, panel & Propiedad, DTO por rol & Indefinida (registro de ventas) \\ \hline
Pago (monto, estado, referencia) & Cobro y conciliación & Financiero & pagos & Dueño del pedido, admin, dueño & Orden, reportes & Token de pasarela nunca sale; idempotencia & Indefinida \\ \hline
Datos de tarjeta & Pagar & Financiero sensible & \textbf{Solo la pasarela} & — & Nunca pasan por Poke Fresh & El cliente paga en la página del proveedor & No se guardan \\ \hline
Credenciales internas, pepper, claves de Mongo y pasarela & Integración & Secreto & Vault & Cada servicio, solo las suyas & Nunca & Políticas por servicio & Hasta su rotación \\ \hline
\end{longtable}
}
\end{landscape}

\subsection{Decisión sobre secretos del cliente}

\textbf{No: Poke Fresh no custodia secretos asociados al cliente.} El cliente paga en la página de la pasarela; ningún dato de tarjeta pasa por el navegador de la tienda, el gateway ni los servicios. Lo que se guarda del pago es el monto, el estado y una referencia de la pasarela, que no permite operar la cuenta de nadie. La contraseña del cliente es una credencial, no un secreto custodiado: se guarda su hash irreversible. El token de la transacción (\texttt{token} en pagos) es un identificador de la operación con alcance limitado a esa transacción; aun así no se devuelve al navegador. La única clave secreta de pago es la del comercio (\texttt{api\_\allowbreak{}key} y \texttt{secreto\_\allowbreak{}firma}), que vive en Vault y nunca llega al navegador.


\subsection{Recorrido de un dato restringido: el RUN}

\begin{enumerate}[nosep,leftmargin=*]
\item \textbf{Captura:} el cliente lo escribe en \texttt{\#/\allowbreak{}registro}. El frontend valida el dígito verificador por comodidad.
\item \textbf{Transporte:} viaja en el cuerpo de \texttt{POST /\allowbreak{}api/\allowbreak{}v1/\allowbreak{}clientes/\allowbreak{}registro} (por HTTPS; el navegador solo habla con el gateway, y con TLS). El gateway no lo registra en logs (solo método, ruta, estado y rol).
\item \textbf{Validación y almacenamiento:} el servicio de clientes vuelve a validar formato, dígito verificador y unicidad (índice único) y lo guarda en \texttt{pf\_\allowbreak{}clientes}, normalizado (\texttt{16459723-\allowbreak{}7}).
\item \textbf{Uso:} al crear un pedido, pedidos copia nombre, RUN, correo y teléfono en el pedido (la boleta debe mostrar los datos del momento de la venta).
\item \textbf{Exposición mínima:} cocina y despacho reciben del pedido solo nombre y teléfono; el RUN lo ven el propio cliente, el administrador y el dueño (boleta y CSV del reporte).
\item \textbf{Eliminación o conservación:} al desactivar al cliente, su ficha queda inactiva y su cuenta no puede ingresar. La copia en pedidos y boletas se conserva justificadamente como respaldo tributario de la venta. Una eliminación de la ficha a petición del titular borraría la ficha y la cuenta, conservando solo los documentos de venta.
\end{enumerate}

\subsection{Sesiones y tokens}

\begin{itemize}[nosep,leftmargin=*]
\item El token es opaco (no es un JWT): no contiene datos y el Auth Service guarda solo su hash.
\item Va en una cookie \texttt{HttpOnly; Secure; SameSite=\allowbreak{}Lax; Path=\allowbreak{}/\allowbreak{}api} (solo viaja por HTTPS). El JavaScript de la página no puede leerlo, así que un script inyectado no lo roba.
\item Expira a los 30 minutos sin uso (cada petición lo renueva) o a las 8 horas.
\item \textbf{Revocación inmediata:} logout, desactivar la cuenta, cambiar el correo, el rol o la clave borran las sesiones; la siguiente petición responde 401.
\item \textbf{Ninguna clave compartida vive en el navegador:} el frontend no tiene API keys ni secretos; todo lo que necesita autenticarse lo hace el gateway.
\end{itemize}

\subsection{Casos de validación ejecutados}

Ejecutados contra el backend implementado con \texttt{pruebas/\allowbreak{}e2e/\allowbreak{}casos.\allowbreak{}py}. Resultados completos y verificables en \texttt{evidencias/\allowbreak{}e2e-\allowbreak{}resultados.\allowbreak{}md} (\texttt{docs/\allowbreak{}evidencias/\allowbreak{}e2e-\allowbreak{}resultados.\allowbreak{}md}) y \texttt{e2e-\allowbreak{}resultados.\allowbreak{}json} (\texttt{docs/\allowbreak{}evidencias/\allowbreak{}e2e-\allowbreak{}resultados.\allowbreak{}json}) (con el \texttt{X-\allowbreak{}Request-\allowbreak{}Id} de cada petición).


{\footnotesize
\setlength{\tabcolsep}{3pt}\renewcommand{\arraystretch}{1.15}
\begin{longtable}{|>{\raggedright\arraybackslash}p{\dimexpr 0.1598\linewidth-2\tabcolsep\relax}|>{\raggedright\arraybackslash}p{\dimexpr 0.2438\linewidth-2\tabcolsep\relax}|>{\raggedright\arraybackslash}p{\dimexpr 0.0940\linewidth-2\tabcolsep\relax}|>{\raggedright\arraybackslash}p{\dimexpr 0.2181\linewidth-2\tabcolsep\relax}|>{\raggedright\arraybackslash}p{\dimexpr 0.1905\linewidth-2\tabcolsep\relax}|>{\raggedright\arraybackslash}p{\dimexpr 0.0940\linewidth-2\tabcolsep\relax}|}
\hline
\textbf{Caso} & \textbf{Entrada} & \textbf{Identidad} & \textbf{Esperado} & \textbf{Obtenido} & \textbf{Lo verifica} \\ \hline
\endfirsthead
\hline
\textbf{Caso} & \textbf{Entrada} & \textbf{Identidad} & \textbf{Esperado} & \textbf{Obtenido} & \textbf{Lo verifica} \\ \hline
\endhead
7.1 Acceso anónimo al catálogo & \texttt{GET /\allowbreak{}catalogo/\allowbreak{}items} & anónimo & 200 sin stock & 200, 29 productos sin stock & Gateway + inventario \\ \hline
7.2 Operación protegida sin sesión & \texttt{GET /\allowbreak{}pedidos} & anónimo & 401 & 401 SIN\_\allowbreak{}SESION & Gateway \\ \hline
7.3 Administración con rol sin permiso & \texttt{DELETE} producto, \texttt{GET /\allowbreak{}clientes}, reportes & cocina, dueño, cliente & 403 & 403 en las cuatro & Gateway + inventario \\ \hline
7.4 Identidad falsificada & \texttt{X-\allowbreak{}Usuario-\allowbreak{}Rol:\allowbreak{} admin} & cliente & 403 & 403 & Gateway \\ \hline
7.5 Pedido ajeno & REST y GraphQL & otro cliente & 404 / null & 404 sin datos; \texttt{resumenPedido:\allowbreak{} null} & pedidos \\ \hline
7.6 Pago alterado desde el frontend & \texttt{monto:\allowbreak{} 1}; notificación sin firma / firma falsa & cliente & monto real; 401 & monto 7.890; 401 y 401; pedido sigue pendiente & pagos \\ \hline
7.7 CSRF & \texttt{Origin} ajeno & cliente & 403 & 403 ORIGEN\_\allowbreak{}NO\_\allowbreak{}PERMITIDO & Gateway \\ \hline
7.8 Cookie y logout & login, logout & cliente & HttpOnly y Secure; 401 tras logout & Cumple & Gateway + Auth \\ \hline
7.9 Acceso directo al backend & desde un contenedor de la red privada, sin la credencial del destino y con \texttt{X-\allowbreak{}Usuario-\allowbreak{}Rol:\allowbreak{} admin} & atacante dentro de la red & 403/401 & pedidos 403, inventario 403, GraphQL 403, Auth 401 & Cada servicio \\ \hline
7.10 Secreto ajeno en Vault & dentro del contenedor de inventario, con su propio token & inventario & 403 & propio 200; pagos 403; pepper 403 & Vault \\ \hline
7.11 Fuerza bruta & 25 logins fallidos & anónimo & bloqueo & 429 desde el intento 6 & Auth + gateway \\ \hline
7.12 Transporte cifrado & GET en claro y por HTTPS & anónimo & sin HTTP en claro; HTTPS con HSTS; certificado verificado con la CA & Cumple & Gateway (TLS) \\ \hline
\end{longtable}
}

\textbf{Qué comprobamos.} Revisamos tres cosas. Primero, que no haya secretos en el JavaScript, el HTML ni las capturas: lo único que se muestra son cuentas de prueba, que el gateway informa solo con \texttt{MODO\_\allowbreak{}DEMO=\allowbreak{}true}. Segundo, que cada rol reciba solo los datos que necesita: un cliente no ve pedidos de otro (caso 7.5) y cocina no ve el RUN ni el correo de quien compra (caso 10.4). Tercero, que aunque la tienda esté abierta a Internet, lo que llega al navegador lo decide el servidor según el rol y no la pantalla (casos 7.2 a 7.4).

